% GENERATED FILE. Edit canonical lessons, then run npm run generate:books.
% canonical-source-hash: fnv1a64:4eef039ef507e201
% canonical-lessons: TE-C339-sainyam, TE-C339-sainikudu, TE-C339-senadhipati, TE-C339-kavatu, TE-C339-dandettu

\chapter{The Army Marches}
\label{ch:te-the-army-marches}
\hlchaptermodality{339}

\begin{chapteropening}
\textbf{By the end of this chapter:} \emph{I can say an army, a soldier, a commander, a march and to march against.}

Complete the last lesson of chapter 339: I can say an army, a soldier, a commander, a march and to march against.
\end{chapteropening}

\section[\texorpdfstring{\te{సైన్యం} (sainyaṁ)}{సైన్యం (sainyaṁ)}]{\te{సైన్యం} (sainyaṁ) — an army}

\label{lesson:TE-C339-sainyam}



\begin{quote}
Before the new one: say the Telugu for the past, then the Telugu for in those days.
\end{quote}

\subsection*{You'll want to know: \te{సైన్యం}}
\textbf{\te{సైన్యం}} — \emph{sainyaṁ} — \textquotedblleft{}an army\textquotedblright{}.

\begin{quote}
\te{శత్రువు} \te{సైన్యం} \te{కోట} \te{దగ్గరకు} \te{వచ్చింది}.
\end{quote}

\begin{quote}
\emph{śatruvu sainyaṁ kōṭa daggaraku vaccindi.}
\end{quote}

\begin{quote}
The enemy army came near the fort.
\end{quote}

Plural: \te{సైన్యాలు}, \emph{sainyālu}.

\begin{cousinweb}
\textbf{\te{సైన్యం}} (\emph{sainyaṁ}) is borrowed from Sanskrit \emph{sainya}, \textquotedblleft{}belonging to an army; an army\textquotedblright{} (CDIAL 13602), built on \emph{senā}, \textquotedblleft{}an army\textquotedblright{}.
\end{cousinweb}

\begin{grammarlens}[title={What you've built}]
One more word for telling what happened.
\end{grammarlens}

\subsection*{Your turn}
\begin{itemize}
\raggedright
  \item \emph{Say it:} \emph{sainyaṁ}
  \item \emph{Say it:} \emph{sainyaṁ}, once more
  \item \emph{Say it:} \emph{śatruvu sainyaṁ kōṭa daggaraku vaccindi.} — The enemy army came near the fort.
  \item \emph{Recall:} say the Telugu for the past, then the Telugu for in those days, then say \emph{sainyaṁ} again
\end{itemize}

\begin{tcolorbox}[breakable,colback=teal!4,colframe=teal!35!black,title={Before you move on}]
What does \te{సైన్యం} mean? (\textquotedblleft{}An army\textquotedblright{}.) Say it once more.
\end{tcolorbox}
\section[\texorpdfstring{\te{సైనికుడు} (sainikuḍu)}{సైనికుడు (sainikuḍu)}]{\te{సైనికుడు} (sainikuḍu) — a soldier}

\label{lesson:TE-C339-sainikudu}



\begin{quote}
Before the new one: say the Telugu for in those days, then the Telugu for an army.
\end{quote}

\subsection*{You'll want to know: \te{సైనికుడు}}
\textbf{\te{సైనికుడు}} — \emph{sainikuḍu} — \textquotedblleft{}a soldier\textquotedblright{}.

\begin{quote}
\te{సైనికుడు} \te{రాత్రంతా} \te{కోటకు} \te{కాపలా} \te{కాశాడు}.
\end{quote}

\begin{quote}
\emph{sainikuḍu rātrantā kōṭaku kāpalā kāśāḍu.}
\end{quote}

\begin{quote}
The soldier guarded the fort all night.
\end{quote}

Plural: \te{సైనికులు}, \emph{sainikulu}.

The -\te{ఇక} ending makes one belonging to something: \te{సైనికుడు}, a man of the army.

\begin{cousinweb}
\textbf{\te{సైనికుడు}} (\emph{sainikuḍu}) is borrowed from Sanskrit \emph{sainika}, \textquotedblleft{}a soldier, a man of the army\textquotedblright{} (Monier-Williams), with the Telugu ending \emph{-uḍu}.
\end{cousinweb}

\begin{grammarlens}[title={What you've built}]
One more word for telling what happened.
\end{grammarlens}

\subsection*{Your turn}
\begin{itemize}
\raggedright
  \item \emph{Say it:} \emph{sainikuḍu}
  \item \emph{Say it:} \emph{sainikuḍu}, once more
  \item \emph{Say it:} \emph{sainikuḍu rātrantā kōṭaku kāpalā kāśāḍu.} — The soldier guarded the fort all night.
  \item \emph{Recall:} say the Telugu for in those days, then the Telugu for an army, then say \emph{sainikuḍu} again
\end{itemize}

\begin{tcolorbox}[breakable,colback=teal!4,colframe=teal!35!black,title={Before you move on}]
What does \te{సైనికుడు} mean? (\textquotedblleft{}A soldier\textquotedblright{}.) Say it once more.
\end{tcolorbox}
\section[\texorpdfstring{\te{సేనాధిపతి} (sēnādhipati)}{సేనాధిపతి (sēnādhipati)}]{\te{సేనాధిపతి} (sēnādhipati) — a commander; a general}

\label{lesson:TE-C339-senadhipati}



\begin{quote}
Before the new one: say the Telugu for an army, then the Telugu for a soldier.
\end{quote}

\subsection*{You'll want to know: \te{సేనాధిపతి}}
\textbf{\te{సేనాధిపతి}} — \emph{sēnādhipati} — \textquotedblleft{}a commander; a general\textquotedblright{}.

\begin{quote}
\te{సేనాధిపతి} \te{సైనికులతో} \te{మాట్లాడాడు}.
\end{quote}

\begin{quote}
\emph{sēnādhipati sainikulatō māṭlāḍāḍu.}
\end{quote}

\begin{quote}
The commander spoke with the soldiers.
\end{quote}

\te{సేన}, an army, and \te{అధిపతి}, a lord or chief.

\begin{cousinweb}
\textbf{\te{సేనాధిపతి}} (\emph{sēnādhipati}) is borrowed from Sanskrit, a compound of \emph{senā}, \textquotedblleft{}an army\textquotedblright{}, and \emph{adhipati}, \textquotedblleft{}a lord, master\textquotedblright{} (CDIAL 255).
\end{cousinweb}

\begin{grammarlens}[title={What you've built}]
One more word for telling what happened.
\end{grammarlens}

\subsection*{Your turn}
\begin{itemize}
\raggedright
  \item \emph{Say it:} \emph{sēnādhipati}
  \item \emph{Say it:} \emph{sēnādhipati}, once more
  \item \emph{Say it:} \emph{sēnādhipati sainikulatō māṭlāḍāḍu.} — The commander spoke with the soldiers.
  \item \emph{Recall:} say the Telugu for an army, then the Telugu for a soldier, then say \emph{sēnādhipati} again
\end{itemize}

\begin{tcolorbox}[breakable,colback=teal!4,colframe=teal!35!black,title={Before you move on}]
What does \te{సేనాధిపతి} mean? (\textquotedblleft{}A commander; a general\textquotedblright{}.) Say it once more.
\end{tcolorbox}
\section[\texorpdfstring{\te{కవాతు} (kavātu)}{కవాతు (kavātu)}]{\te{కవాతు} (kavātu) — a march; a parade}

\label{lesson:TE-C339-kavatu}



\begin{quote}
Before the new one: say the Telugu for a soldier, then the Telugu for a commander.
\end{quote}

\subsection*{You'll want to know: \te{కవాతు}}
\textbf{\te{కవాతు}} — \emph{kavātu} — \textquotedblleft{}a march; a parade\textquotedblright{}.

\begin{quote}
\te{సైనికులు} \te{మైదానంలో} \te{కవాతు} \te{చేశారు}.
\end{quote}

\begin{quote}
\emph{sainikulu maidānaṁlō kavātu cēśāru.}
\end{quote}

\begin{quote}
The soldiers marched on the parade ground.
\end{quote}

\te{కవాతు} \te{చేయు}, to march in step.

\begin{cousinweb}
\textbf{\te{కవాతు}} (\emph{kavātu}) is the Telugu word for a drill march; its origin is not settled.
\end{cousinweb}

\begin{grammarlens}[title={What you've built}]
One more word for telling what happened.
\end{grammarlens}

\subsection*{Your turn}
\begin{itemize}
\raggedright
  \item \emph{Say it:} \emph{kavātu}
  \item \emph{Say it:} \emph{kavātu}, once more
  \item \emph{Say it:} \emph{sainikulu maidānaṁlō kavātu cēśāru.} — The soldiers marched on the parade ground.
  \item \emph{Recall:} say the Telugu for a soldier, then the Telugu for a commander, then say \emph{kavātu} again
\end{itemize}

\begin{tcolorbox}[breakable,colback=teal!4,colframe=teal!35!black,title={Before you move on}]
What does \te{కవాతు} mean? (\textquotedblleft{}A march; a parade\textquotedblright{}.) Say it once more.
\end{tcolorbox}
\section[\texorpdfstring{\te{దండెత్తు} (daṇḍettu)}{దండెత్తు (daṇḍettu)}]{\te{దండెత్తు} (daṇḍettu) — to march against; to invade}

\label{lesson:TE-C339-dandettu}



\begin{quote}
Before the new one: say the Telugu for a commander, then the Telugu for a march.
\end{quote}

\subsection*{You'll want to know: \te{దండెత్తు}}
\textbf{\te{దండెత్తు}} — \emph{daṇḍettu} — \textquotedblleft{}to march against; to invade\textquotedblright{}.

\begin{quote}
\te{ఆ} \te{రాజు} \te{పొరుగు} \te{రాజ్యం} \te{మీద} \te{దండెత్తాడు}.
\end{quote}

\begin{quote}
\emph{ā rāju porugu rājyaṁ mīda daṇḍettāḍu.}
\end{quote}

\begin{quote}
That king marched against the neighbouring kingdom.
\end{quote}

Past (I …): \te{దండెత్తాను}, \emph{daṇḍettānu}.

The target takes \te{మీద}: X \te{మీద} \te{దండెత్తు}.

\begin{cousinweb}
\textbf{\te{దండెత్తు}} (\emph{daṇḍettu}) is a native compound: \emph{daṇḍu}, \textquotedblleft{}an army, a host\textquotedblright{} (DEDR 3055), and \emph{ettu}, \textquotedblleft{}to lift, raise; to set out\textquotedblright{} (DEDR 796): to raise an army against someone.
\end{cousinweb}

\begin{grammarlens}[title={What you've built}]
One more word for telling what happened.
\end{grammarlens}

\subsection*{Your turn}
\begin{itemize}
\raggedright
  \item \emph{Say it:} \emph{daṇḍettu}
  \item \emph{Say it:} \emph{daṇḍettu}, once more
  \item \emph{Say it:} \emph{ā rāju porugu rājyaṁ mīda daṇḍettāḍu.} — That king marched against the neighbouring kingdom.
  \item \emph{Recall:} say the Telugu for a commander, then the Telugu for a march, then say \emph{daṇḍettu} again
\end{itemize}

\begin{tcolorbox}[breakable,colback=teal!4,colframe=teal!35!black,title={Before you move on}]
What does \te{దండెత్తు} mean? (\textquotedblleft{}To march against; to invade\textquotedblright{}.) Say it once more.
\end{tcolorbox}
